\IfFileExists{stacks-project.cls}{%
\documentclass{stacks-project}
}{%
\documentclass{amsart}
}

% For dealing with references we use the comment environment
\usepackage{verbatim}
\newenvironment{reference}{\comment}{\endcomment}
%\newenvironment{reference}{}{}
\newenvironment{slogan}{\comment}{\endcomment}
\newenvironment{history}{\comment}{\endcomment}

% For commutative diagrams we use Xy-pic
\usepackage[all]{xy}

% We use 2cell for 2-commutative diagrams.
\xyoption{2cell}
\UseAllTwocells

% We use multicol for the list of chapters between chapters
\usepackage{multicol}

% This is generally recommended for better output
\usepackage{lmodern}
\usepackage[T1]{fontenc}

% For cross-file-references
\usepackage{xr-hyper}

% Package for hypertext links:
\usepackage{hyperref}

% For any local file, say "hello.tex" you want to link to please
% use \externaldocument[hello-]{hello}
\externaldocument[introduction-]{introduction}
\externaldocument[conventions-]{conventions}
\externaldocument[sets-]{sets}
\externaldocument[categories-]{categories}
\externaldocument[topology-]{topology}
\externaldocument[sheaves-]{sheaves}
\externaldocument[sites-]{sites}
\externaldocument[stacks-]{stacks}
\externaldocument[fields-]{fields}
\externaldocument[algebra-]{algebra}
\externaldocument[brauer-]{brauer}
\externaldocument[homology-]{homology}
\externaldocument[derived-]{derived}
\externaldocument[simplicial-]{simplicial}
\externaldocument[more-algebra-]{more-algebra}
\externaldocument[smoothing-]{smoothing}
\externaldocument[modules-]{modules}
\externaldocument[sites-modules-]{sites-modules}
\externaldocument[injectives-]{injectives}
\externaldocument[cohomology-]{cohomology}
\externaldocument[sites-cohomology-]{sites-cohomology}
\externaldocument[dga-]{dga}
\externaldocument[dpa-]{dpa}
\externaldocument[sdga-]{sdga}
\externaldocument[hypercovering-]{hypercovering}
\externaldocument[schemes-]{schemes}
\externaldocument[constructions-]{constructions}
\externaldocument[properties-]{properties}
\externaldocument[morphisms-]{morphisms}
\externaldocument[coherent-]{coherent}
\externaldocument[divisors-]{divisors}
\externaldocument[limits-]{limits}
\externaldocument[varieties-]{varieties}
\externaldocument[topologies-]{topologies}
\externaldocument[descent-]{descent}
\externaldocument[perfect-]{perfect}
\externaldocument[more-morphisms-]{more-morphisms}
\externaldocument[flat-]{flat}
\externaldocument[groupoids-]{groupoids}
\externaldocument[more-groupoids-]{more-groupoids}
\externaldocument[etale-]{etale}
\externaldocument[chow-]{chow}
\externaldocument[intersection-]{intersection}
\externaldocument[pic-]{pic}
\externaldocument[weil-]{weil}
\externaldocument[adequate-]{adequate}
\externaldocument[dualizing-]{dualizing}
\externaldocument[duality-]{duality}
\externaldocument[discriminant-]{discriminant}
\externaldocument[derham-]{derham}
\externaldocument[local-cohomology-]{local-cohomology}
\externaldocument[algebraization-]{algebraization}
\externaldocument[curves-]{curves}
\externaldocument[resolve-]{resolve}
\externaldocument[models-]{models}
\externaldocument[functors-]{functors}
\externaldocument[equiv-]{equiv}
\externaldocument[pione-]{pione}
\externaldocument[etale-cohomology-]{etale-cohomology}
\externaldocument[proetale-]{proetale}
\externaldocument[relative-cycles-]{relative-cycles}
\externaldocument[more-etale-]{more-etale}
\externaldocument[trace-]{trace}
\externaldocument[crystalline-]{crystalline}
\externaldocument[spaces-]{spaces}
\externaldocument[spaces-properties-]{spaces-properties}
\externaldocument[spaces-morphisms-]{spaces-morphisms}
\externaldocument[decent-spaces-]{decent-spaces}
\externaldocument[spaces-cohomology-]{spaces-cohomology}
\externaldocument[spaces-limits-]{spaces-limits}
\externaldocument[spaces-divisors-]{spaces-divisors}
\externaldocument[spaces-over-fields-]{spaces-over-fields}
\externaldocument[spaces-topologies-]{spaces-topologies}
\externaldocument[spaces-descent-]{spaces-descent}
\externaldocument[spaces-perfect-]{spaces-perfect}
\externaldocument[spaces-more-morphisms-]{spaces-more-morphisms}
\externaldocument[spaces-flat-]{spaces-flat}
\externaldocument[spaces-groupoids-]{spaces-groupoids}
\externaldocument[spaces-more-groupoids-]{spaces-more-groupoids}
\externaldocument[bootstrap-]{bootstrap}
\externaldocument[spaces-pushouts-]{spaces-pushouts}
\externaldocument[spaces-chow-]{spaces-chow}
\externaldocument[groupoids-quotients-]{groupoids-quotients}
\externaldocument[spaces-more-cohomology-]{spaces-more-cohomology}
\externaldocument[spaces-simplicial-]{spaces-simplicial}
\externaldocument[spaces-duality-]{spaces-duality}
\externaldocument[formal-spaces-]{formal-spaces}
\externaldocument[restricted-]{restricted}
\externaldocument[spaces-resolve-]{spaces-resolve}
\externaldocument[formal-defos-]{formal-defos}
\externaldocument[defos-]{defos}
\externaldocument[cotangent-]{cotangent}
\externaldocument[examples-defos-]{examples-defos}
\externaldocument[algebraic-]{algebraic}
\externaldocument[examples-stacks-]{examples-stacks}
\externaldocument[stacks-sheaves-]{stacks-sheaves}
\externaldocument[criteria-]{criteria}
\externaldocument[artin-]{artin}
\externaldocument[quot-]{quot}
\externaldocument[stacks-properties-]{stacks-properties}
\externaldocument[stacks-morphisms-]{stacks-morphisms}
\externaldocument[stacks-limits-]{stacks-limits}
\externaldocument[stacks-cohomology-]{stacks-cohomology}
\externaldocument[stacks-perfect-]{stacks-perfect}
\externaldocument[stacks-introduction-]{stacks-introduction}
\externaldocument[stacks-more-morphisms-]{stacks-more-morphisms}
\externaldocument[stacks-geometry-]{stacks-geometry}
\externaldocument[moduli-]{moduli}
\externaldocument[moduli-curves-]{moduli-curves}
\externaldocument[examples-]{examples}
\externaldocument[exercises-]{exercises}
\externaldocument[guide-]{guide}
\externaldocument[desirables-]{desirables}
\externaldocument[coding-]{coding}
\externaldocument[obsolete-]{obsolete}
\externaldocument[fdl-]{fdl}
\externaldocument[index-]{index}

% Theorem environments.
%
\theoremstyle{plain}
\newtheorem{theorem}[subsection]{Theorem}
\newtheorem{proposition}[subsection]{Proposition}
\newtheorem{lemma}[subsection]{Lemma}

\theoremstyle{definition}
\newtheorem{definition}[subsection]{Definition}
\newtheorem{example}[subsection]{Example}
\newtheorem{exercise}[subsection]{Exercise}
\newtheorem{situation}[subsection]{Situation}

\theoremstyle{remark}
\newtheorem{remark}[subsection]{Remark}
\newtheorem{remarks}[subsection]{Remarks}

\numberwithin{equation}{subsection}

% Macros
%
\def\lim{\mathop{\mathrm{lim}}\nolimits}
\def\colim{\mathop{\mathrm{colim}}\nolimits}
\def\Spec{\mathop{\mathrm{Spec}}}
\def\Hom{\mathop{\mathrm{Hom}}\nolimits}
\def\Ext{\mathop{\mathrm{Ext}}\nolimits}
\def\SheafHom{\mathop{\mathcal{H}\!\mathit{om}}\nolimits}
\def\SheafExt{\mathop{\mathcal{E}\!\mathit{xt}}\nolimits}
\def\Sch{\mathit{Sch}}
\def\Mor{\mathop{\mathrm{Mor}}\nolimits}
\def\Ob{\mathop{\mathrm{Ob}}\nolimits}
\def\Sh{\mathop{\mathit{Sh}}\nolimits}
\def\NL{\mathop{N\!L}\nolimits}
\def\CH{\mathop{\mathrm{CH}}\nolimits}
\def\proetale{{pro\text{-}\acute{e}tale}}
\def\etale{{\acute{e}tale}}
\def\QCoh{\mathit{QCoh}}
\def\Ker{\mathop{\mathrm{Ker}}}
\def\Im{\mathop{\mathrm{Im}}}
\def\Coker{\mathop{\mathrm{Coker}}}
\def\Coim{\mathop{\mathrm{Coim}}}

% Boxtimes
%
\DeclareMathSymbol{\boxtimes}{\mathbin}{AMSa}{"02}

%
% Macros for moduli stacks/spaces
%
\def\QCohstack{\mathcal{QC}\!\mathit{oh}}
\def\Cohstack{\mathcal{C}\!\mathit{oh}}
\def\Spacesstack{\mathcal{S}\!\mathit{paces}}
\def\Quotfunctor{\mathrm{Quot}}
\def\Hilbfunctor{\mathrm{Hilb}}
\def\Curvesstack{\mathcal{C}\!\mathit{urves}}
\def\Polarizedstack{\mathcal{P}\!\mathit{olarized}}
\def\Complexesstack{\mathcal{C}\!\mathit{omplexes}}
% \Pic is the operator that assigns to X its picard group, usage \Pic(X)
% \Picardstack_{X/B} denotes the Picard stack of X over B
% \Picardfunctor_{X/B} denotes the Picard functor of X over B
\def\Pic{\mathop{\mathrm{Pic}}\nolimits}
\def\Picardstack{\mathcal{P}\!\mathit{ic}}
\def\Picardfunctor{\mathrm{Pic}}
\def\Deformationcategory{\mathcal{D}\!\mathit{ef}}
\usepackage{cleveref}
\crefname{originallemma}{original lemma}{original lemmas}
\Crefname{originallemma}{Original lemma}{Original lemmas}
\makeatletter
\newlabel{lemma-ameliorate-morphism-categories-fibred-groupoids}{{06N7}{source}{}{sourcepointer-1}{}}
\newlabel{lemma-ameliorate-morphism-categories-fibred-groupoids@cref}{{[originallemma][0][]06N7}{[1][0][]source}{}{}{}}
\makeatother
\crefname{originallemma}{original lemma}{original lemmas}
\Crefname{originallemma}{Original lemma}{Original lemmas}
\makeatletter
\newlabel{lemma-equivalence-2-fibre-product}{{02XB}{source}{}{sourcepointer-2}{}}
\newlabel{lemma-equivalence-2-fibre-product@cref}{{[originallemma][0][]02XB}{[1][0][]source}{}{}{}}
\makeatother
\crefname{originallemma}{original lemma}{original lemmas}
\Crefname{originallemma}{Original lemma}{Original lemmas}
\makeatletter
\newlabel{lemma-equivalence-fibred-categories}{{003Z}{source}{}{sourcepointer-3}{}}
\newlabel{lemma-equivalence-fibred-categories@cref}{{[originallemma][0][]003Z}{[1][0][]source}{}{}{}}
\makeatother
\title{Uniqueness of the amelioration of a 1-morphism of categories fibred in groupoids}
\author{The Stacks Project Authors}
\date{stacks-project:categories.tex:local-snapshot-20261012; GNU FDL 1.2 or later}
\begin{document}
\maketitle
\noindent\textit{Editorial note.} Counted span is the author's own proof environment, source lines 7793-7844, following the statement at 7771-7791.\par\medskip
\noindent\textit{Editorial note.} The identification of the comparison object uses the in-chapter construction referenced explicitly; all 2-isomorphism data is displayed in-span.\par\medskip
\noindent\textit{Allowed standard prerequisites (editorial).}\begin{itemize}
\item 2-fibre product of categories over C. Lemma lemma-2-product-categories-over-C of this chapter; the identification of the ameliorated category with X x\_\{F,Y,id\} Y is given in the proof of lemma-ameliorate-morphism-categories-fibred-groupoids at 7725-7728. The comparison object over Y to which both X' and X'' are related.
\item Equivalences and 2-fibre products. Lemma lemma-equivalence-2-fibre-product of this chapter, invoked at proof lines 7813-7814. c is an equivalence because b is.
\item Equivalences of categories fibred in groupoids. Lemma lemma-equivalence-fibred-categories of this chapter, invoked at proof lines 7815-7817. Promotes the equivalence over C to an equivalence in the 2-category of categories fibred in groupoids over Y.
\end{itemize}\medskip
\ifdefined\hypertarget\hypertarget{sourcepointer-1}{}\fi
\noindent\textit{Source reference.} Original lemma 06N7: sources/stacks-categories/pinned-original-source.tex, source lines 7707--7707. Navigation pointer only: the counted proof cites this statement by \textbackslash{}ref as a standing input; the deduction the counted proof has to make is carried out inside the counted span, and no part of that deduction is delegated to this pointer.\par\medskip
\ifdefined\hypertarget\hypertarget{sourcepointer-2}{}\fi
\noindent\textit{Source reference.} Original lemma 02XB: sources/stacks-categories/pinned-original-source.tex, source lines 5933--5933. Navigation pointer only: the counted proof cites this statement by \textbackslash{}ref as a standing input; the deduction the counted proof has to make is carried out inside the counted span, and no part of that deduction is delegated to this pointer.\par\medskip
\ifdefined\hypertarget\hypertarget{sourcepointer-3}{}\fi
\noindent\textit{Source reference.} Original lemma 003Z: sources/stacks-categories/pinned-original-source.tex, source lines 7488--7488. Navigation pointer only: the counted proof cites this statement by \textbackslash{}ref as a standing input; the deduction the counted proof has to make is carried out inside the counted span, and no part of that deduction is delegated to this pointer.\par\medskip
\begin{lemma}
\label{lemma-amelioration-unique}
Let $\mathcal{C}$ be a category. Let $F : \mathcal{X} \to \mathcal{Y}$
be a $1$-morphism of categories fibred in groupoids over $\mathcal{C}$.
Assume we have a $2$-commutative diagram
$$
\xymatrix{
\mathcal{X}' \ar[rd]_f &
\mathcal{X} \ar[l]^a \ar[d]^F \ar[r]_b &
\mathcal{X}'' \ar[ld]^g \\
& \mathcal{Y}
}
$$
where $a$ and $b$ are equivalences of categories over $\mathcal{C}$
and $f$ and $g$ are categories fibred in groupoids. Then there exists
an equivalence $h : \mathcal{X}'' \to \mathcal{X}'$ of categories over
$\mathcal{Y}$ such that $h \circ b$ is $2$-isomorphic to $a$ as $1$-morphisms
of categories over $\mathcal{C}$. If the diagram above actually commutes, then
we can arrange it so that $h \circ b$ is $2$-isomorphic to $a$ as
$1$-morphisms of categories over $\mathcal{Y}$.
\end{lemma}

\begin{proof}
We will show that both $\mathcal{X}'$ and $\mathcal{X}''$ over $\mathcal{Y}$
are equivalent to the category fibred in groupoids
$\mathcal{X} \times_{F, \mathcal{Y}, \text{id}} \mathcal{Y}$
over $\mathcal{Y}$, see proof of
Lemma \ref{lemma-ameliorate-morphism-categories-fibred-groupoids}.
Choose a quasi-inverse $b^{-1} : \mathcal{X}'' \to \mathcal{X}$ in the
$2$-category of categories over $\mathcal{C}$.
Since the right triangle of the diagram is $2$-commutative we see that
$$
\xymatrix{
\mathcal{X} \ar[d]_F & \mathcal{X}'' \ar[l]^{b^{-1}} \ar[d]^g \\
\mathcal{Y} & \mathcal{Y} \ar[l]
}
$$
is $2$-commutative. Hence we obtain a $1$-morphism
$c : \mathcal{X}'' \to
\mathcal{X} \times_{F, \mathcal{Y}, \text{id}} \mathcal{Y}$
by the universal property of the $2$-fibre product. Moreover $c$
is a morphism of categories over $\mathcal{Y}$ (!) and an equivalence
(by the assumption that $b$ is an equivalence, see
Lemma \ref{lemma-equivalence-2-fibre-product}).
Hence $c$ is an equivalence in the $2$-category of categories fibred
in groupoids over $\mathcal{Y}$ by
Lemma \ref{lemma-equivalence-fibred-categories}.

\medskip\noindent
We still have to construct a $2$-isomorphism between $c \circ b$ and
the functor $d : \mathcal{X} \to
\mathcal{X} \times_{F, \mathcal{Y}, \text{id}} \mathcal{Y}$,
$x \mapsto (p(x), x, F(x), \text{id}_{F(x)})$
constructed in the proof of
Lemma \ref{lemma-ameliorate-morphism-categories-fibred-groupoids}.
Let $\alpha : F \to g \circ b$ and $\beta : b^{-1} \circ b \to \text{id}$
be $2$-isomorphisms between $1$-morphisms of categories over $\mathcal{C}$.
Note that $c \circ b$ is given by the rule
$$
x \mapsto (p(x), b^{-1}(b(x)), g(b(x)), \alpha_x \circ F(\beta_x))
$$
on objects. Then we see that
$$
(\beta_x, \alpha_x) :
(p(x), x, F(x), \text{id}_{F(x)})
\longrightarrow
(p(x), b^{-1}(b(x)), g(b(x)), \alpha_x \circ F(\beta_x))
$$
is a functorial isomorphism which gives our $2$-morphism
$d \to b \circ c$. Finally, if the diagram commutes then
$\alpha_x$ is the identity for all $x$ and we see that this
$2$-morphism is a $2$-morphism in the $2$-category of categories
over $\mathcal{Y}$.
\end{proof}

\end{document}
